% Pengaturan ukuran teks dan bentuk halaman dua sisi
\documentclass[12pt]{book}

% Pengaturan ukuran halaman dan margin
\usepackage[a4paper,top=30mm,left=30mm,right=20mm,bottom=25mm]{geometry}

% Pengaturan ukuran spasi
\usepackage[singlespacing]{setspace}

% Pengaturan caption untuk tabel
\usepackage{caption}

% Judul dokumen
\title{Proposal Tugas Akhir ITS}
\author{Musk, Elon Reeve}

% Pengaturan detail pada file PDF
\usepackage[pdfauthor={\@author},bookmarksnumbered,pdfborder={0 0 0}]{hyperref}


% Pengaturan ukuran indentasi
\setlength{\parindent}{2em}

% Package lainnya
\usepackage{changepage}
\usepackage{etoolbox} % Mengubah fungsi default

% Pengaturan jenis karakter
\usepackage[utf8]{inputenc}

\usepackage[style=ieee, backend=biber]{biblatex}
\usepackage{enumitem} % Pembuatan list
\usepackage{lipsum} % Pembuatan template kalimat
\usepackage{graphicx} % Input gambar
\usepackage{longtable} % Pembuatan tabel
\usepackage[table,xcdraw]{xcolor} % Pewarnaan tabel
\usepackage{eso-pic} % Untuk menggunakan background image di halaman
\usepackage{txfonts} % Font times
\usepackage{changepage} % Pembuatan teks kolom
\usepackage{multicol} % Pembuatan kolom ganda
\usepackage{multirow} % Pembuatan baris ganda
\usepackage{tabularx} % Untuk mengatur kolom, seperti grid pada CSS
\usepackage{wrapfig}
\usepackage{float}

% Pengaturan format daftar isi, daftar gambar, dan daftar tabel
\usepackage[titles]{tocloft}
\setlength{\cftsecindent}{2em}
\setlength{\cftsubsecindent}{2em}
\setlength{\cftbeforechapskip}{1.5ex}
\setlength{\cftbeforesecskip}{1.5ex}
\setlength{\cftbeforetoctitleskip}{0cm}
\setlength{\cftbeforeloftitleskip}{0cm}
\setlength{\cftbeforelottitleskip}{0cm}
\renewcommand{\cfttoctitlefont}{\hfill\Large\bfseries} % command untuk membuat heading bold dan besar
\renewcommand{\cftaftertoctitle}{\hfill}
\renewcommand{\cftloftitlefont}{\hfill\Large\bfseries}
\renewcommand{\cftafterloftitle}{\hfill}
\renewcommand{\cftlottitlefont}{\hfill\Large\bfseries}
\renewcommand{\cftafterlottitle}{\hfill}

% Definisi untuk "Hati ini sengaja dikosongkan"
\patchcmd{\cleardoublepage}{\hbox{}}{
  \thispagestyle{empty}
  \vspace*{\fill}
  \begin{center}\textit{[Halaman ini sengaja dikosongkan]}\end{center}
  \vfill}{}{}

  % Pengaturan penomoran halaman
\usepackage{fancyhdr}
\fancyhf{}
\renewcommand{\headrulewidth}{0pt}
\pagestyle{fancy}
\fancyfoot[C,CO]{\thepage}
\patchcmd{\chapter}{plain}{fancy}{}{}
\patchcmd{\chapter}{empty}{plain}{}{}

% Pengaturan format judul bab
\usepackage{titlesec}
\renewcommand{\thesection}{\thechapter.\arabic{section}}
\titleformat{\chapter}[hang]{\centering\bfseries\large}{BAB\ \arabic{chapter}\ }{0ex}{\vspace{0ex}\centering}
\titleformat*{\section}{\large\bfseries}
\titleformat*{\subsection}{\normalsize\bfseries}
\titlespacing{\chapter}{0ex}{0ex}{4ex}
\titlespacing{\section}{0ex}{1ex}{0ex}
\titlespacing{\subsection}{0ex}{0.5ex}{0ex}
\titlespacing{\subsubsection}{0ex}{0.5ex}{0ex}
\setcounter{secnumdepth}{3} % Untuk memberi penomoran pada \subsubsection

\counterwithin{figure}{chapter}
\counterwithin{table}{chapter}

% Mengganti figure dan table menjadi gambar dan tabel
\renewcommand{\figurename}{Gambar}
\renewcommand{\tablename}{Tabel}

% Tambahkan format tanda hubung yang benar di sini
\hyphenation{
  ro-ket
  me-ngem-bang-kan
  per-hi-tu-ngan
}

% Menambahkan resource daftar pustaka
\addbibresource{pustaka/pustaka.bib}

% Isi keseluruhan dokumen
\begin{document}
  % Nomor halaman pembuka dimulai dari sini
  \pagenumbering{roman}

  % Atur ulang penomoran halaman
  \setcounter{page}{1}

  % Sampul Bahasa Indonesia
  \newcommand\covercontents{sampul/konten-id.tex}
  \AddToShipoutPictureBG*{
  \AtPageLowerLeft{
    % Ubah nilai berikut jika posisi horizontal background tidak sesuai
    \hspace{-3.25mm}

    % Ubah nilai berikut jika posisi vertikal background tidak sesuai
    \raisebox{0mm}{
      \includegraphics[width=\paperwidth,height=\paperheight]{sampul/gambar/sampul-luar-tipis.png}
    }
  }
}

% Menyembunyikan nomor halaman
\thispagestyle{empty}

% Pengaturan margin untuk menyesuaikan konten sampul
\newgeometry{
  top=65mm,
  left=30mm,
  right=30mm,
  bottom=20mm
}

\begin{flushleft}

  % Pemilihan font sans serif
  \sffamily

  % Pemilihan font bold
  \fontseries{bx}
  \selectfont
  \begin{spacing}{1.5}
    \input{\covercontents}
  \end{spacing}

\end{flushleft}

\restoregeometry


  % Lembar pengesahan
  \chapter*{LEMBAR PENGESAHAN}

% Menyembunyikan nomor halaman
\thispagestyle{empty}

\begin{center}
  % Ubah kalimat berikut dengan judul tugas akhir
  \textbf{KALKULASI ENERGI PADA ROKET LUAR ANGKASA BERBASIS \emph{ANTI-GRAVITASI}}
\end{center}

\begingroup
% Pemilihan font ukuran small
\small

\begin{center}
  % Ubah kalimat berikut dengan pernyataan untuk lembar pengesahan
  \textbf{PROPOSAL TUGAS AKHIR} \\
  Diajukan untuk memenuhi salah satu syarat memperoleh gelar
  Sarjana Teknik pada
  Program Studi S-1 Teknik Dirgantara \\
  Departemen Teknik Dirgantara \\
  Fakultas Teknik Dirgantara \\
  Institut Teknologi Sepuluh Nopember
\end{center}

\begin{center}
  % Ubah kalimat berikut dengan nama dan NRP mahasiswa
  Oleh: \textbf{Elon Reeve Musk} \\
  NRP. 0123 20 4000 0001
\end{center}

\begin{center}
  Disetujui Oleh:
\end{center}

\vspace{10ex}

\begingroup
% Menghilangkan padding
\setlength{\tabcolsep}{0pt}

\noindent
\begin{tabularx}{\textwidth}{X c}
  % Ubah kalimat-kalimat berikut dengan nama dan NIP dosen pembimbing pertama
  Nikola Tesla, S.T., M.T.      &                 \\
  NIP: 18560710 194301 1 001    & (Pembimbing)    \\
                                &                 \\
                                &                 \\
                                &                 \\
  % Ubah kalimat-kalimat berikut dengan nama dan NIP dosen pembimbing kedua
  Wernher von Braun, S.T., M.T. &                 \\
  NIP: 19230323 197706 1 001    & (Ko-Pembimbing) \\
\end{tabularx}
\endgroup

\vspace{\fill}

\begin{center}
  Mengetahui,\\
  % Ubah kalimat berikut dengan nama departemen
  Kepala Departemen Teknik Dirgantara FTD-ITS\\
  \vspace{10ex}
  % Ubah kalimat berikut dengan jabatan kepala departemen
  \underline{Nikola Tesla, S.T., M.T. }\\
  NIP 18560710 194301 1 001\\
  \vspace{10ex}
  % Ubah text dibawah menjadi tempat dan tanggal
  \textbf{SURABAYA} \\
  \textbf{Mei, 2077}
\end{center}
\endgroup

  \cleardoublepage

  % Abstrak
  % Halaman Kata Pengantar
\chapter*{KATA PENGANTAR}
\addcontentsline{toc}{chapter}{KATA PENGANTAR}

Segala puji dan syukur kami panjatkan ke hadirat Tuhan Yang Maha Esa atas limpahan rahmat dan karunia-Nya, sehingga kami dapat menyelesaikan Laporan Proyek Telematika dengan judul \textbf{Praktika App} ini dengan baik. Dalam proses penyusunan laporan ini, kami ingin menyampaikan rasa terima kasih yang sebesar-besarnya kepada pihak-pihak yang telah memberikan dukungan, baik secara langsung maupun tidak langsung, sehingga laporan ini dapat terselesaikan dengan lancar, yaitu:

\begin{itemize}[leftmargin=2cm]
    \item Tuhan Yang Maha Esa atas segala kelancaran, keberkahan, dan bimbingan yang diberikan selama pelaksanaan proyek ini.
    \item Bapak Dr. Supeno Mardi Susiki Nugroho, S.T., M.T., selaku Kepala Departemen Teknik Komputer FTEIC-ITS, atas dukungan dan arahannya.
    \item Bapak Ahmad Zaini, S.T., M.T., Dr. Arief Kurniawan, S.T., M.T., dan Prof. Dr. I Ketut Eddy Purnama, S.T., M.T., selaku para Dosen Pengampu Proyek Telematika Departemen Teknik Komputer FTEIC-ITS, atas bimbingan dan ilmunya selama proses penyusunan laporan ini.
    \item Rekan-rekan kelompok proyek telematika, yang telah bekerja sama dengan penuh dedikasi selama pelaksanaan proyek ini.
    \item Rekan satu tim dan semua pihak lainnya yang tidak dapat kami sebutkan satu per satu namun telah memberikan kontribusi yang sangat berarti.
\end{itemize}

Kami menyadari bahwa laporan ini masih jauh dari kesempurnaan. Oleh karena itu, kami memohon maaf jika terdapat kekurangan atau hal-hal yang kurang berkenan selama pelaksanaan proyek maupun dalam penyusunan laporan ini. Kami berharap, laporan ini dapat memberikan manfaat, baik bagi pembaca maupun bagi pengembangan keilmuan di bidang telematika di masa yang akan datang. Demikian, semoga laporan ini dapat menjadi referensi dan inspirasi bagi para pembaca.

\vspace{1cm}

\begin{flushright}
Surabaya, Desember 2024 \\
\vspace{4cm}
Penulis
\end{flushright}

  \cleardoublepage

  \chapter*{ABSTRACT}
\begin{center}
  \large
  \textbf{\emph{ANTI-GRAVITY} BASED ENERGY CALCULATION ON OUTER SPACE ROCKETS}
\end{center}
% Menyembunyikan nomor halaman
\thispagestyle{empty}

\begin{flushleft}
  \setlength{\tabcolsep}{0pt}
  \bfseries
  \begin{tabular}{lc@{\hspace{6pt}}l}
  Student Name / NRP&: &Elon Reeve Musk / 0123204000001\\
  Department&: &Aerospace Engineering FTD - ITS\\
  Advisor&: &1. Nikola Tesla, S.T., M.T.\\
  & & 2. Wernher von Braun, S.T., M.T.\\
  \end{tabular}
  \vspace{4ex}
\end{flushleft}
\textbf{Abstract}

% Isi Abstrak
The abstract must consist between two hundred to three hundred words. \lipsum[1]

\vspace{2ex}
\noindent
\textbf{Keywords: \emph{Rocket, Anti-gravity, Meong}}
  \cleardoublepage

  \begin{spacing}{1.5}
    % Daftar isi
    \renewcommand*\contentsname{DAFTAR ISI}
    \addcontentsline{toc}{chapter}{\contentsname}
    \tableofcontents
    \cleardoublepage

    % Daftar gambar
    \renewcommand*\listfigurename{DAFTAR GAMBAR}
    \addcontentsline{toc}{chapter}{\listfigurename}
    \listoffigures
    \cleardoublepage

    % Daftar tabel
    \renewcommand*\listtablename{DAFTAR TABEL}
    \addcontentsline{toc}{chapter}{\listtablename}
    \listoftables
    \cleardoublepage
  \end{spacing}

  % Nomor halaman isi dimulai dari sini
  \pagenumbering{arabic}

  % Konten pendahuluan
  \chapter{PENDAHULUAN}

\section{Latar Belakang}

% Ubah paragraf-paragraf berikut sesuai dengan latar belakang dari tugas akhir
Pesatnya perkembangan roket yang merupakan \lipsum[2]

\lipsum[3]

\section{Rumusan Masalah}

% Ubah paragraf berikut sesuai dengan rumusan masalah dari tugas akhir
Berdasarkan hal yang telah dipaparkan di latar belakang, \lipsum[4]

\section{Batasan Masalah atau Ruang Lingkup}

\lipsum[6]

\section{Tujuan}

% Ubah paragraf berikut sesuai dengan tujuan penelitian dari tugas akhir
Tujuan dari penelitian ini adalah \lipsum[7][1-14]

\section{Manfaat}

% Ubah paragraf berikut sesuai dengan tujuan penelitian dari tugas akhir
Manfaat dari penelitian ini adalah \lipsum[8][1-14]

  \cleardoublepage

  % Konten tinjauan pustaka
  \chapter{TINJAUAN PUSTAKA}

% Ubah konten-konten berikut sesuai dengan isi dari tinjauan pustaka
\section{Hasil penelitian/perancangan terdahulu}
\lipsum[3]

\section{Teori/Konsep Dasar}

\subsection{Hukum Newton}

% Contoh penggunaan referensi dari pustaka
Newton pernah merumuskan \parencite{Newton1687} bahwa \lipsum[8]
% Contoh penggunaan referensi dari persamaan
Kemudian menjadi persamaan seperti pada persamaan \ref{eq:FirstLaw}.

% Contoh pembuatan persamaan
\begin{equation}
  % Label referensi dari persamaan yang dibuat
  \label{eq:FirstLaw}
  % Baris kode persamaan yang dibuat
  \sum \mathbf{F} = 0\; \Leftrightarrow\; \frac{\mathrm{d} \mathbf{v} }{\mathrm{d}t} = 0.
\end{equation}

\lipsum[9]

\subsection{Anti Gravitasi}

\lipsum[10]

  \cleardoublepage

  % Konten metodologi
  \chapter{KESIMPULAN DAN SARAN}


\section{Kesimpulan}

Berdasarkan hasil pengembangan dan implementasi sistem, dapat disimpulkan bahwa sistem backend yang dibuat untuk aplikasi PraktikaApp, PraktikaSSO, dan PraktikaTelelab berjalan dengan baik. Sistem ini telah berhasil di-deploy secara lokal dan dilengkapi dengan berbagai fitur keamanan yang mendukung, seperti \emph{OpenID Auth} untuk otentikasi pengguna. Selain itu, sistem backend ini dibangun dengan pendekatan \emph{microservice}, yang memungkinkan pengelolaan dan pengembangan aplikasi secara modular dan terpisah sesuai dengan fungsinya.

Selain itu, untuk memastikan kemudahan deployment dan skalabilitas, aplikasi ini dibungkus menggunakan Docker dan Kubernetes. Dengan demikian, sistem dapat dengan mudah di-deploy pada berbagai lingkungan dan mendukung proses pengelolaan kontainer yang efisien. Fitur keamanan lain yang diterapkan adalah Cross-Origin Resource Sharing (\emph{CORS}), yang memberikan perlindungan terhadap akses tidak sah dari sumber yang berbeda, memastikan aplikasi tetap aman dan terkontrol.

\section{Saran}

Sebagai pengembangan lebih lanjut, ada beberapa saran yang dapat dipertimbangkan untuk meningkatkan kualitas dan fungsionalitas sistem, antara lain:

\begin{itemize}
    \item Pengembangan aplikasi \textbf{Praktika App} dapat diperluas untuk mencakup berbagai komponen selain rangkaian digital, sehingga dapat mencakup berbagai aspek yang lebih luas dalam dunia pendidikan dan teknologi.
    \item Melakukan \emph{deployment} aplikasi di lingkungan \emph{production} yang saat ini belum dilakukan, agar aplikasi dapat digunakan oleh pengguna secara lebih luas dan memberikan feedback yang lebih komprehensif terkait performa dan kestabilan aplikasi di dunia nyata.
    \item Menambahkan fitur baru yang dapat memperkaya pengalaman pengguna, seperti analisis data penggunan aplikasi atau integrasi dengan layanan pihak ketiga untuk meningkatkan fungsionalitas sistem.
\end{itemize}

  \cleardoublepage

  % Konten lainnya
  \chapter{HASIL YANG DIHARAPKAN}

\section{Hasil yang Diharapkan dari Penelitian}

Dari penelitian yang akan dilakukan, diharapkan \lipsum[15]

\section{Hasil Pendahuluan}

Sampai saat ini, kami telah \lipsum[16]

  \cleardoublepage

  \chapter{JADWAL PENELITIAN}

% Ubah tabel berikut sesuai dengan isi dari rencana kerja
\newcommand{\w}{}
\newcommand{\G}{\cellcolor{gray}}
\begin{table}[H]
  \captionof{table}{Tabel timeline}
  \label{tbl:timeline}
  \begin{tabular}{|p{3.5cm}|c|c|c|c|c|c|c|c|c|c|c|c|c|c|c|c|}

    \hline
    \multirow{2}{*}{Kegiatan} & \multicolumn{16}{|c|}{Minggu}                                                                       \\
    \cline{2-17}              &
    1                         & 2                             & 3  & 4  & 5  & 6  & 7  & 8  & 9  & 10 & 11 & 12 & 13 & 14 & 15 & 16 \\
    \hline

    % Gunakan \G untuk mengisi sel dan \w untuk mengosongkan sel
    Pengambilan data          &
    \G                        & \G                            & \G & \G & \w & \w & \w & \w & \w & \w & \w & \w & \w & \w & \w & \w \\
    \hline

    Pengolahan data           &
    \w                        & \w                            & \w & \w & \G & \G & \G & \G & \w & \w & \w & \w & \w & \w & \w & \w \\
    \hline

    Analisa data              &
    \w                        & \w                            & \w & \w & \w & \w & \w & \w & \G & \G & \G & \G & \w & \w & \w & \w \\
    \hline

    Evaluasi penelitian       &
    \w                        & \w                            & \w & \w & \w & \w & \w & \w & \w & \w & \w & \w & \G & \G & \G & \G \\
    \hline
  \end{tabular}
\end{table}

Pada \emph{timeline} yang tertera di Tabel \ref{tbl:timeline} \lipsum[10]

  \cleardoublepage

  % Daftar pustaka
  \chapter*{DAFTAR PUSTAKA}
  \addcontentsline{toc}{chapter}{DAFTAR PUSTAKA}
  \renewcommand\refname{}
  \vspace{2ex}
  \renewcommand{\bibname}{}
  \begingroup
    \def\chapter*#1{}
    \printbibliography
  \endgroup


\end{document}
